\documentclass[10pt,a4paper]{amsart}
\usepackage{amsmath,amssymb,amsthm,graphicx}
% Retained minimal preamble for the Stacks Project chapter flat.tex (More on Flatness).
% Verbatim excerpt of the chapter preamble: Xy-pic setup (native preamble lines 14-18) and the native
% theorem declarations with the mathematical macros (native preamble lines 154-217). Copied mechanically;
% nothing added, removed or redefined. The wrapper class and the amsmath/amssymb/amsthm/graphicx base
% packages are supplied by the assembler, so the class/externaldocument part of the native preamble is omitted.
% For commutative diagrams we use Xy-pic
\usepackage[all]{xy}

% We use 2cell for 2-commutative diagrams.
\xyoption{2cell}
% Theorem environments.
%
\theoremstyle{plain}
\newtheorem{theorem}[subsection]{Theorem}
\newtheorem{proposition}[subsection]{Proposition}
\newtheorem{lemma}[subsection]{Lemma}

\theoremstyle{definition}
\newtheorem{definition}[subsection]{Definition}
\newtheorem{example}[subsection]{Example}
\newtheorem{exercise}[subsection]{Exercise}
\newtheorem{situation}[subsection]{Situation}

\theoremstyle{remark}
\newtheorem{remark}[subsection]{Remark}
\newtheorem{remarks}[subsection]{Remarks}

\numberwithin{equation}{subsection}

% Macros
%
\def\lim{\mathop{\mathrm{lim}}\nolimits}
\def\colim{\mathop{\mathrm{colim}}\nolimits}
\def\Spec{\mathop{\mathrm{Spec}}}
\def\Hom{\mathop{\mathrm{Hom}}\nolimits}
\def\Ext{\mathop{\mathrm{Ext}}\nolimits}
\def\SheafHom{\mathop{\mathcal{H}\!\mathit{om}}\nolimits}
\def\SheafExt{\mathop{\mathcal{E}\!\mathit{xt}}\nolimits}
\def\Sch{\mathit{Sch}}
\def\Mor{\mathop{\mathrm{Mor}}\nolimits}
\def\Ob{\mathop{\mathrm{Ob}}\nolimits}
\def\Sh{\mathop{\mathit{Sh}}\nolimits}
\def\NL{\mathop{N\!L}\nolimits}
\def\CH{\mathop{\mathrm{CH}}\nolimits}
\def\proetale{{pro\text{-}\acute{e}tale}}
\def\etale{{\acute{e}tale}}
\def\QCoh{\mathit{QCoh}}
\def\Ker{\mathop{\mathrm{Ker}}}
\def\Im{\mathop{\mathrm{Im}}}
\def\Coker{\mathop{\mathrm{Coker}}}
\def\Coim{\mathop{\mathrm{Coim}}}

% Boxtimes
%
\DeclareMathSymbol{\boxtimes}{\mathbin}{AMSa}{"02}

%
% Macros for moduli stacks/spaces
%
\def\QCohstack{\mathcal{QC}\!\mathit{oh}}
\def\Cohstack{\mathcal{C}\!\mathit{oh}}
\def\Spacesstack{\mathcal{S}\!\mathit{paces}}
\def\Quotfunctor{\mathrm{Quot}}
\def\Hilbfunctor{\mathrm{Hilb}}
\def\Curvesstack{\mathcal{C}\!\mathit{urves}}
\def\Polarizedstack{\mathcal{P}\!\mathit{olarized}}
\def\Complexesstack{\mathcal{C}\!\mathit{omplexes}}
% \Pic is the operator that assigns to X its picard group, usage \Pic(X)
% \Picardstack_{X/B} denotes the Picard stack of X over B
% \Picardfunctor_{X/B} denotes the Picard functor of X over B
\def\Pic{\mathop{\mathrm{Pic}}\nolimits}
\def\Picardstack{\mathcal{P}\!\mathit{ic}}
\def\Picardfunctor{\mathrm{Pic}}
\def\Deformationcategory{\mathcal{D}\!\mathit{ef}}

% --- appended 2026-10-12 (Chief): the project own definitions of the reference/slogan/history
% display environments, copied mechanically from the shared Stacks preamble lines 8-12
% (\usepackage{verbatim} plus the three \newenvironment lines, which the chapter preamble keeps
% verbatim from the project preamble). Nothing added, removed or redefined: these environments
% carry no mathematical content and the project itself typesets them as comments. Without these
% three lines a counted statement whose own text uses them fails with Environment slogan undefined.
\usepackage{verbatim}
\newenvironment{reference}{\comment}{\endcomment}
\newenvironment{slogan}{\comment}{\endcomment}
\newenvironment{history}{\comment}{\endcomment}
\usepackage{cleveref}
\crefname{originaldefinition}{original definition}{original definitions}
\Crefname{originaldefinition}{Original definition}{Original definitions}
\makeatletter
\newlabel{topologies-definition-big-small-ph}{{0DBL}{source}{}{sourcepointer-1}{}}
\newlabel{topologies-definition-big-small-ph@cref}{{[originaldefinition][0][]0DBL}{[1][0][]source}{}{}{}}
\makeatother
\crefname{originallemma}{original lemma}{original lemmas}
\Crefname{originallemma}{Original lemma}{Original lemmas}
\makeatletter
\newlabel{topologies-lemma-characterize-sheaf}{{0DBQ}{source}{}{sourcepointer-2}{}}
\newlabel{topologies-lemma-characterize-sheaf@cref}{{[originallemma][0][]0DBQ}{[1][0][]source}{}{}{}}
\makeatother
\crefname{originallemma}{original lemma}{original lemmas}
\Crefname{originallemma}{Original lemma}{Original lemmas}
\makeatletter
\newlabel{morphisms-lemma-seminormalization}{{0EUS}{source}{}{sourcepointer-3}{}}
\newlabel{morphisms-lemma-seminormalization@cref}{{[originallemma][0][]0EUS}{[1][0][]source}{}{}{}}
\makeatother
\crefname{originallemma}{original lemma}{original lemmas}
\Crefname{originallemma}{Original lemma}{Original lemmas}
\makeatletter
\newlabel{morphisms-lemma-scheme-theoretic-image-is-proper}{{0AH6}{source}{}{sourcepointer-4}{}}
\newlabel{morphisms-lemma-scheme-theoretic-image-is-proper@cref}{{[originallemma][0][]0AH6}{[1][0][]source}{}{}{}}
\makeatother
\crefname{originallemma}{original lemma}{original lemmas}
\Crefname{originallemma}{Original lemma}{Original lemmas}
\makeatletter
\newlabel{morphisms-lemma-finite-proper}{{01WN}{source}{}{sourcepointer-5}{}}
\newlabel{morphisms-lemma-finite-proper@cref}{{[originallemma][0][]01WN}{[1][0][]source}{}{}{}}
\makeatother
\crefname{originallemma}{original lemma}{original lemmas}
\Crefname{originallemma}{Original lemma}{Original lemmas}
\makeatletter
\newlabel{morphisms-lemma-universally-injective}{{01S4}{source}{}{sourcepointer-6}{}}
\newlabel{morphisms-lemma-universally-injective@cref}{{[originallemma][0][]01S4}{[1][0][]source}{}{}{}}
\makeatother
\crefname{originallemma}{original lemma}{original lemmas}
\Crefname{originallemma}{Original lemma}{Original lemmas}
\makeatletter
\newlabel{morphisms-lemma-universal-homeomorphism}{{04DF}{source}{}{sourcepointer-7}{}}
\newlabel{morphisms-lemma-universal-homeomorphism@cref}{{[originallemma][0][]04DF}{[1][0][]source}{}{}{}}
\makeatother
\title{The absolutely weakly normal structure sheaf is a sheaf for the ph topology}
\author{The Stacks Project Authors}
\date{stacks-project:flat.tex:local-snapshot-20261011; GNU FDL 1.2 or later}
\begin{document}
\maketitle
\noindent\textit{Editorial note.} All delivered mathematical text is the author's own, copied from the pinned Stacks Project chapter file flat.tex; no statement or proof step was rewritten, reordered or removed.\par\medskip
\noindent\textit{Editorial note.} The preamble is the chapter's own retained minimal preamble (sources/stacks-flat/retained-preamble.tex), a verbatim excerpt of the native Xy-pic setup and of the native theorem declarations and mathematical macros, followed by the project own definitions of the reference, slogan and history display environments (verbatim package plus three newenvironment lines) that the retained file carries at its end and that the project itself typesets as comments. No macro was invented or redefined, and the wrapper class and the amsmath/amssymb/amsthm/graphicx base packages are supplied by the assembler.\par\medskip
\noindent\textit{Editorial note.} The counted statement and proof contain no cross reference, no citation and no included file, so the delivered body is self contained.\par\medskip
\noindent\textit{Editorial note.} Canonical location of the counted proposition: Stacks tag 0EVT (sources/stacks-flat/pinned-original-source.tex line 11991).\par\medskip
\noindent\textit{Allowed standard prerequisites (editorial).}\begin{itemize}
\item The characterisation of ph sheaves by two conditions. Topologies, Lemma labelled topologies-lemma-characterize-sheaf, cited at proof lines 12001-12003. Reduces the sheaf property to the Zariski and the equalizer conditions.
\item Commutation of the absolute weak normalization with open coverings. Morphisms, Lemma labelled morphisms-lemma-seminormalization, cited at proof lines 12005-12008. Settles condition (1).
\item Scheme theoretic images of proper morphisms and finiteness of proper morphisms with affine target. Morphisms, Lemmas labelled morphisms-lemma-scheme-theoretic-image-is-proper and morphisms-lemma-finite-proper, cited at proof lines 12020-12027. Produces the closed subscheme Z of A\textasciicircum{}1\_X proper and finite over X.
\item Universal injectivity and universal homeomorphisms. Morphisms, Lemmas labelled morphisms-lemma-universally-injective and morphisms-lemma-universal-homeomorphism, cited at proof lines 12060-12068. Converts the pointwise injectivity proved by the field argument into the universal homeomorphism applied at the end.
\end{itemize}\medskip
\ifdefined\hypertarget\hypertarget{sourcepointer-1}{}\fi
\noindent\textit{Source reference.} Original definition 0DBL: sources/stacks-chow/tags.txt, source lines 16701--16701. Cross-chapter standard prerequisite located in the project tag list staged in the local snapshot; invoked at proof line None of the counted proof. It is retained only as a navigation pointer; every load-bearing step is inside the delivered span.\par\medskip
\ifdefined\hypertarget\hypertarget{sourcepointer-2}{}\fi
\noindent\textit{Source reference.} Original lemma 0DBQ: sources/stacks-chow/tags.txt, source lines 16705--16705. Cross-chapter standard prerequisite located in the project tag list staged in the local snapshot; invoked at proof line 12003 of the counted proof. It is retained only as a navigation pointer; every load-bearing step is inside the delivered span.\par\medskip
\ifdefined\hypertarget\hypertarget{sourcepointer-3}{}\fi
\noindent\textit{Source reference.} Original lemma 0EUS: sources/stacks-chow/tags.txt, source lines 18562--18562. Cross-chapter standard prerequisite located in the project tag list staged in the local snapshot; invoked at proof line 12005 of the counted proof. It is retained only as a navigation pointer; every load-bearing step is inside the delivered span.\par\medskip
\ifdefined\hypertarget\hypertarget{sourcepointer-4}{}\fi
\noindent\textit{Source reference.} Original lemma 0AH6: sources/stacks-chow/tags.txt, source lines 13219--13219. Cross-chapter standard prerequisite located in the project tag list staged in the local snapshot; invoked at proof line 12023 of the counted proof. It is retained only as a navigation pointer; every load-bearing step is inside the delivered span.\par\medskip
\ifdefined\hypertarget\hypertarget{sourcepointer-5}{}\fi
\noindent\textit{Source reference.} Original lemma 01WN: sources/stacks-chow/tags.txt, source lines 2488--2488. Cross-chapter standard prerequisite located in the project tag list staged in the local snapshot; invoked at proof line 12025 of the counted proof. It is retained only as a navigation pointer; every load-bearing step is inside the delivered span.\par\medskip
\ifdefined\hypertarget\hypertarget{sourcepointer-6}{}\fi
\noindent\textit{Source reference.} Original lemma 01S4: sources/stacks-chow/tags.txt, source lines 2325--2325. Cross-chapter standard prerequisite located in the project tag list staged in the local snapshot; invoked at proof line 12061 of the counted proof. It is retained only as a navigation pointer; every load-bearing step is inside the delivered span.\par\medskip
\ifdefined\hypertarget\hypertarget{sourcepointer-7}{}\fi
\noindent\textit{Source reference.} Original lemma 04DF: sources/stacks-chow/tags.txt, source lines 5725--5725. Cross-chapter standard prerequisite located in the project tag list staged in the local snapshot; invoked at proof line 12063 of the counted proof. It is retained only as a navigation pointer; every load-bearing step is inside the delivered span.\par\medskip
\begin{lemma}
\label{lemma-weak-normalization-ph-sheaf}
Let $(\Sch/S)_{ph}$ be a site as in
Topologies, Definition \ref{topologies-definition-big-small-ph}.
The rule
$$
X \longmapsto \Gamma(X^{awn}, \mathcal{O}_{X^{awn}})
$$
is a sheaf on $(\Sch/S)_{ph}$.
\end{lemma}

\begin{proof}
To prove $\mathcal{F}$ is a sheaf, let's check conditions
(1) and (2) of Topologies, Lemma \ref{topologies-lemma-characterize-sheaf}.
Condition (1) holds because formation of $X^{awn}$ commutes with
open coverings, see Morphisms, Lemma \ref{morphisms-lemma-seminormalization}
and its proof.

\medskip\noindent
Let $\pi : Y \to X$ be a surjective proper morphism. We have to show
that the equalizer of the two maps
$$
\Gamma(Y^{awn}, \mathcal{O}_{Y^{awn}}) \to
\Gamma((Y \times_X Y)^{awn}, \mathcal{O}_{(Y \times_X Y)^{awn}})
$$
is equal to $\Gamma(X^{awn}, \mathcal{O}_{X^{awn}})$. Let $f$ be
an element of this equalizer. Then we consider the morphism
$$
f : Y^{awn} \longrightarrow \mathbf{A}^1_X
$$
Since $Y^{awn} \to X$ is universally closed, the scheme theoretic
image $Z$ of $f$ is a closed subscheme of $\mathbf{A}^1_X$ proper over $X$
and $f : Y^{awn} \to Z$ is surjective.
See Morphisms, Lemma \ref{morphisms-lemma-scheme-theoretic-image-is-proper}.
Thus $Z \to X$ is finite
(Morphisms, Lemma \ref{morphisms-lemma-finite-proper})
and surjective.

\medskip\noindent
Let $k$ be a field and let $z_1, z_2 : \Spec(k) \to Z$ be two
morphisms equalized by $Z \to X$. We claim that $z_1 = z_2$.
It suffices to show the images $\lambda_i = z_i^*f \in k$ agree
(as the structure sheaf of $Z$ is generated by $f$
over the structure sheaf of $X$). To see this we
choose a field extension
$K/k$ and morphisms $y_1, y_2 : \Spec(K) \to Y^{awn}$ such that
$z_i \circ (\Spec(K) \to \Spec(k)) = f \circ y_i$. This is possible
by the surjectivity of the map $Y^{awn} \to Z$. Choose an algebraically
closed extension $\Omega/k$ of very large cardinality.
For any $k$-algebra maps $\sigma_i : K \to \Omega$
we obtain
$$
\Spec(\Omega)
\xrightarrow{\sigma_1, \sigma_2}
\Spec(K \otimes_k K)
\xrightarrow{y_1, y_2}
Y^{awn} \times_X Y^{awn}
$$
Since the canonical morphism
$(Y \times_X Y)^{awn} \to Y^{awn} \times_X Y^{awn}$
is a universal homeomorphism and since $\Omega$ is algebraically closed,
we can lift the composition above uniquely to a morphism
$\Spec(\Omega) \to (Y \times_X Y)^{awn}$. Since $f$ is in the equalizer
above, this proves that $\sigma_1(\lambda_1) = \sigma_2(\lambda_2)$.
An easy lemma about field extensions shows that this implies
$\lambda_1 = \lambda_2$; details omitted.

\medskip\noindent
We conclude that $Z \to X$ is universally injective, i.e.,
$Z \to X$ is injective on points and induces
purely inseparated residue field extensions
(Morphisms, Lemma \ref{morphisms-lemma-universally-injective}).
All in all we conclude that $Z \to X$ is a universal homeomorphism, see
Morphisms, Lemma \ref{morphisms-lemma-universal-homeomorphism}.

\medskip\noindent
Let $g : X^{awn} \to Z$ be the map obtained from the universal property
of $X^{awn}$. Then $Y^{awn} \to X^{awn} \to Z$ and $f : Y^{awn} \to Z$
are two morphisms over $X$. By the universal property of
$Y^{awn} \to Y$ the two corresponding morphisms
$Y^{awn} \to Y \times_X Z$ over $Y$ have to be equal. This implies
that $g \circ \pi^{wan} = f$ as morphisms into $\mathbf{A}^1_X$
and we conclude that $g \in \Gamma(X^{awn}, \mathcal{O}_{X^{awn}})$
is the element we were looking for.
\end{proof}

\end{document}
